\documentclass[11pt]{article}

\usepackage[margin=1in]{geometry}
\usepackage{amsmath, amssymb}
\usepackage{booktabs}
\usepackage{xcolor}
\usepackage{enumitem}
\usepackage{hyperref}
\usepackage{float}

\definecolor{notecol}{rgb}{0.15,0.35,0.6}
\definecolor{todocol}{rgb}{0.7,0.3,0.0}

% One decision / measurement per entry, appended at the end in date order:
%   \entry{ID. Title}{YYYY-MM-DD}{anchor}{Topic}
% IDs: D decision, M measurement, O open, T tooling. Fields in the order listed below.
% A new date starts a new day heading; \topicindex lists the entries by topic.
\makeatletter
\newcommand{\nb@lastdate}{}
\def\nb@topics{}
\newcommand{\entry}[4]{%
  \def\nb@date{#2}%
  \ifx\nb@date\nb@lastdate\else
    \section*{#2}\addcontentsline{toc}{section}{#2}%
    \global\let\nb@lastdate\nb@date
  \fi
  \subsection*{#1 \hfill \small\normalfont\textsc{#4}}%
  \addcontentsline{toc}{subsection}{#1}\label{#3}%
  \nb@addtopic{#4}{#1}{#3}}
\newcommand{\nb@addtopic}[3]{%
  \@ifundefined{nb@topic@#1}{%
    \expandafter\gdef\csname nb@topic@#1\endcsname{}%
    \g@addto@macro\nb@topics{\nb@printtopic{#1}}}{}%
  \expandafter\g@addto@macro\csname nb@topic@#1\endcsname{\item \hyperref[#3]{#2}}}
\newcommand{\nb@printtopic}[1]{\item[#1]%
  \begin{itemize}[nosep]\csname nb@topic@#1\endcsname\end{itemize}}
\newcommand{\topicindex}{\ifx\nb@topics\@empty\else
  \section*{Entries by topic}\addcontentsline{toc}{section}{Entries by topic}%
  \begin{description}[style=nextline]\nb@topics\end{description}\fi}
\makeatother
\newcommand{\field}[1]{\par\smallskip\noindent\textbf{#1.}\ }
\newcommand{\verify}[1]{\textcolor{todocol}{[to verify: #1]}}
\newcommand{\file}[1]{\texttt{\detokenize{#1}}}

\title{Lab Notebook: ML-Based Floor-Planning for Small-Scale SoC}
\author{Daniel Vigna}
\date{Started 2026-10-14}

\begin{document}
\maketitle

\section*{How this notebook works}
Every design decision and every measurement gets one entry, written when it happens and appended in date order.
Each day starts a new heading, and every entry names the topic it touches (see \emph{Entries by topic} at the end).
The fields of an entry: Question, Options, Method, Result, Decision, Revisit if.
Claims that are reasoning but not yet measured are marked \verify{like this}.
Scripts live in \file{scripts/}, raw data in \file{data/}, next to this file (\file{~/papers/notebook/ml-soc} on the VM).

\paragraph{Project.}
A learning-based floorplanner for a small SoC: a reference block set modeled after an RP2350-class
microcontroller die (CPU cores, GPU, PCIe, USB PHY, SRAM, flash control, ADC, programmable I/O, analog and clock
circuitry), with blocks as arbitrary polygons rather than rectangles and per-block power density informed by ASAP7.
A PyTorch model proposes placements, written to an intermediate format and rendered by a custom viewer.
Measured block data from the SRAM compute-in-memory macro (the \file{sram-cim} notebook) feeds the block generator.

\tableofcontents

% =====================================================================
% Entries go here, in date order:
% \entry{D1. Title}{YYYY-MM-DD}{d1}{Topic}
% \field{Question} ...
% \field{Options} ...
% \field{Method} ...
% \field{Result} ...
% \field{Decision} ...
% \field{Revisit if} ...

% =====================================================================
% \section*{Open items}
% \begin{itemize}
%   \item ...
% \end{itemize}

\topicindex

\end{document}
